\chapter{무게가 다른 화살표}

지금까지 화살표는 모두 똑같았어요.
하지만 실제 기록은 그렇지 않아요.
스티커를 한 번 보낸 친구와 열 번 보낸 친구는 다르지요.
어제 보낸 것과 일 년 전에 보낸 것도 달라요.

\section{굵은 화살표}

화살표마다 \textbf{무게}를 적어 둘 수 있어요.
하준이가 서연이에게 스티커를 아홉 번, 도윤이에게 한 번 보냈다면
서연이 쪽 화살표에는 9, 도윤이 쪽 화살표에는 1을 적어요.

쪽지 놀이도 조금 바꿔요.
하준이가 “화살표를 따라가” 카드를 뽑으면, 열 번 가운데 아홉 번은 서연이에게, 한 번은 도윤이에게 쪽지를 넘겨요.
물 붓기라면 하준이의 물 가운데 90\%가 서연이에게, 10\%가 도윤이에게 가요.
무거운 화살표로 물이 더 많이 흐르는 거예요.

\section{많이 보내도 한 번}

그럼 한 번에 스티커를 100장 보낸 것은 1장 보낸 것보다 100배 무거울까요?
이 장에서 소개할 방법은 그렇게 세지 않아요.
보낸 양이 아주 적지만 않으면, 1장을 보냈든 100장을 보냈든 거의 똑같이 “한 번”으로 세요.
양이 늘면 무게도 조금 늘기는 하지만, 금방 1에 닿고 그 뒤로는 더 커지지 않거든요.

이렇게 하면 엄청나게 큰 송금 하나가 지도를 혼자 차지하지 못해요.
화살표의 무게는 “얼마나 많이 보냈나”보다 “몇 번 보냈나”에 가까워져요.

\section{오래된 일은 조금 가볍게}

일 년 전에 친구에게 스티커를 많이 준 일이 있다고 해 봐요.
그 일은 지금도 중요할까요?
조금은 중요하지만, 어제 일만큼은 아니겠지요.
사람의 기억도 시간이 지나면 흐려지고요.

그래서 화살표의 무게에 \textbf{신선도}를 곱해요.
신선도는 오늘 일어난 일이면 거의 1에 가깝고, 시간이 지날수록 줄어들어요.
제 논문에서 쓴 신선도는 이렇게 생겼어요.

\begin{figure}[h]
\centering
\begin{tikzpicture}
\begin{axis}[
  width=0.92\linewidth, height=5.2cm,
  xmin=0, xmax=400, ymin=0, ymax=1.05,
  xlabel={얼마나 오래된 일인가 (일)}, ylabel={신선도},
  xtick={0,90,180,270,365}, ytick={0,0.5,1},
  grid=major, grid style={gray!25},
  tick label style={font=\scriptsize}, label style={font=\small},
]
\addplot[domain=0:400, samples=120, line width=1.4pt, color=sendcol] {1/(1+exp(0.01*(x-180)))};
\addplot[only marks, mark=*, mark size=2pt, color=inkblue] coordinates {(0,0.858) (90,0.711) (180,0.5) (365,0.136)};
\node[font=\scriptsize, anchor=south west] at (axis cs:2,0.86) {오늘: 0.86};
\node[font=\scriptsize, anchor=south west] at (axis cs:92,0.72) {석 달 전: 0.71};
\node[font=\scriptsize, anchor=north east] at (axis cs:178,0.49) {반년 전: 0.5};
\node[font=\scriptsize, anchor=north east] at (axis cs:362,0.11) {일 년 전: 0.14};
\end{axis}
\end{tikzpicture}
\caption{시간이 지날수록 줄어드는 신선도. 반년(180일) 전 일은 무게가 절반이 돼요.}
\end{figure}

반년 전 일은 무게가 절반, 일 년 전 일은 7분의 1쯤 돼요.
그래서 오래전에 송금을 자주 했더라도 요즘 아무 일도 없으면 점수가 조금씩 내려가요.

\section{선생님은 누구에게 쪽지를 줄까}

6장의 쪽지 놀이에서 “선생님께” 카드가 나오면, 선생님은 쪽지를 아무에게나 똑같이 나누어 주었어요.
이 장의 방법에서는 선생님이 조금 달라요.
요즘 여러 친구에게 스티커를 보낸 친구에게 쪽지를 더 자주 줘요.
보낸 상대가 많을수록, 그리고 최근에 보냈을수록 더 자주 받아요.
한 번도 보낸 적이 없는 친구는 선생님에게서 쪽지를 받지 못해요.
그런 친구는 다른 친구들이 보내 준 화살표로만 점수를 얻어요.

\section{이미 나와 있는 방법: 송금 랭크}

이제 지금까지 만든 것을 모두 합쳐 볼게요.

\begin{enumerate}
\item 송금 줄로 지도를 그린다.
\item 송금 하나하나를 “한 번”쯤으로 세고, 신선도를 곱해서 그 화살표의 무게에 더한다.
\item 선생님은 요즘 활발하게 보낸 친구에게 쪽지를 더 자주 준다.
\item 그 지도에서 쪽지 돌리기 놀이(페이지랭크)를 한다.
\end{enumerate}

이 방법에는 이름이 있어요.
2023년에 열린 한 국제 학회에서 도, 도, 응우옌이라는 세 연구자가 블록체인 지갑의 평판을 재려고 발표한 방법이에요.
우리말로 옮기면 “적응형 가중 페이지랭크”이고, 논문에서는 영어 이름의 머리글자를 따서 AWP라고 불러요.
이 책에서는 송금 지도에서 하는 페이지랭크라는 뜻으로 \textbf{송금 랭크}라고 부를게요.
송금 기록으로 지갑의 평판을 재는, 이미 나와 있는 방법이지요.

제 논문에서 송금 랭크는 출발점이에요.
송금 기록으로 평판을 재는 방법이 이미 있으니, 허락 기록으로 만든 점수가 그것과 무엇이 다른지 살펴볼 수 있어요.
그렇다고 두 점수가 시합을 하는 것은 아니에요.
둘은 서로 다른 기록을 보고 서로 다른 것을 재거든요.
그래서 이 책에서는 “누가 이겼나”를 묻지 않고, 여러 점수를 한 줄에 함께 세워 놓고 점수마다 무엇을 잘 맞히는지 볼 거예요.
허락 기록으로 만든 점수는 다음 장에서 소개할게요.

\section{원래 모습 그대로}

이미 나와 있는 방법을 제 마음대로 바꾸면, 그건 더 이상 그 방법이라고 할 수 없겠지요.
그래서 송금 랭크는 원래 논문에 적힌 대로 만들었어요.
송금을 “한 번”쯤으로 세는 것도, 선생님이 활발한 친구에게 쪽지를 더 자주 주는 것도 원래 논문에 있는 규칙이에요.

다만 원래 논문에는 적혀 있지 않은 숫자가 몇 개 있었어요.
신선도가 얼마나 빨리 줄어드는지, 보낸 양을 얼마나 빨리 “한 번”으로 줄이는지 같은 숫자예요.
이 숫자들은 제가 정했고, 논문에 그대로 밝혀 두었어요.

\begin{deeper}[송금 랭크의 식]
신선도는 로지스틱 함수 $\sigma(\Delta t) = 1/\bigl(1+e^{k(\Delta t - t_0)}\bigr)$이고,
논문에서는 $k=0.01$(하루당), $t_0=180$일을 썼어요.
$\Delta t$는 점수를 매기는 날을 기준으로 그 송금이 며칠 전에 일어났는지예요.

송금 하나의 무게는 $\sigma(\Delta t)\times V(z)$이고,
$V(z) = 2/(1+e^{-bz}) - 1$은 보낸 양 $z$를 0과 1 사이로 줄이는 함수예요.
원래 논문에는 $b$ 값이 없어서 $b=1$로 정했어요.
토큰의 양은 아주 작은 단위의 정수로 기록되기 때문에, 그 단위로 10 이상이면 $V(z)$는 1과 0.0001도 차이 나지 않아요.
그래서 송금 한 번이 오래된 정도만큼 줄어든 채 “한 번”으로 세어져요.

되돌아가는 몫(선생님 몫)은 각 지갑의 \emph{활발함}에 비례해서 나누어요.
활발함은 그 지갑이 보낸 상대마다, 그 상대에게 가장 최근에 보낸 송금의 신선도를 더한 값이에요.
\end{deeper}

\begin{learned}
\begin{itemize}
\item 화살표에 무게를 주면, 무거운 화살표로 쪽지가 더 자주 가요.
\item 송금 랭크는 송금 하나를 양과 거의 상관없이 “한 번”으로 세요.
\item 오래된 일은 신선도를 곱해서 가볍게 쳐요. 반년 전 일은 절반이에요.
\item 송금 랭크에서는 선생님이 요즘 활발하게 보낸 친구에게 쪽지를 더 자주 줘요.
\item 송금 랭크는 송금 기록으로 평판을 재는, 이미 나와 있는 방법이에요. 원래 논문에 적힌 대로 만들었어요.
\end{itemize}
\end{learned}
